\documentclass[../../../include/open-logic-section]{subfiles}

\begin{document}
\olfileid{sth}{card-arithmetic}{simp}

\olsection{যোগ ও গুণের সরলীকরণ}

দেখা যাবে, ট্রান্সফাইনাইট অঙ্কবাচক সংখ্যার যোগ ও গুণ
\emph{অত্যন্ত} সহজ। কারণ অঙ্কবাচক সংখ্যাগুলি বিশেষ
কিছু ক্রমসংখ্যা; তাই তারা সুক্রমবিন্যস্ত, এবং সেই
কাঠামো কাজে লাগানো যায়। অবশ্য এটি দেখানো
ততটা সহজ নয়। শুরুতে একটি কৌশলী সংজ্ঞা দরকার:

\begin{defn}
ক্রমসংখ্যার যুগলের উপর একটি \emph{মানক ক্রম}
$\canonord$ সংজ্ঞায়িত করি।
$\tuple{\alpha_1,\alpha_2}\canonord
\tuple{\beta_1,\beta_2}$ তখনই, যখন নিচের
যেকোনো একটি শর্ত পূরণ হয়:
\begin{enumerate}
	\item $\max(\alpha_1, \alpha_2) < \max(\beta_1, \beta_2)$; অথবা
	\item $\max(\alpha_1, \alpha_2) = \max(\beta_1, \beta_2)$ এবং
	$\alpha_1 < \beta_1$; অথবা
	\item $\max(\alpha_1, \alpha_2) = \max(\beta_1, \beta_2)$ এবং
	$\alpha_1 = \beta_1$ এবং $\alpha_2 < \beta_2$।
\end{enumerate}
\end{defn}

\begin{lem}
যেকোনো ক্রমসংখ্যা $\alpha$-র জন্য
$\tuple{\alpha\times\alpha,\canonord}$ একটি সুক্রমবিন্যাস।
\end{lem}

\begin{proof}
$\alpha\times\alpha$-র যেকোনো দুটি ভিন্ন যুগল
$\canonord$-এ তুলনাযোগ্য। কারণ
$\tuple{\alpha_1,\alpha_2}$ ও
$\tuple{\beta_1,\beta_2}$-এর কোনোটিই অন্যটির
চেয়ে $\canonord$-এ ছোট না হলে
$\max(\alpha_1,\alpha_2)=\max(\beta_1,\beta_2)$,
$\alpha_1=\beta_1$ এবং $\alpha_2=\beta_2$।
অতএব $\tuple{\alpha_1,\alpha_2}=
\tuple{\beta_1,\beta_2}$।

সুক্রমবিন্যাস দেখাতে অশূন্য $X\subseteq\alpha\times\alpha$
ধরি। $\alpha$ ক্রমসংখ্যা বলে
$\Setabs{\max(\gamma_1,\gamma_2)}{\tuple{\gamma_1,
\gamma_2}\in X}$-এর ক্ষুদ্রতম সদস্য কোনো
$\delta$। এবার
$\Setabs{\tuple{\gamma_1,\gamma_2}\in X}{\max(\gamma_1,\gamma_2)=
\delta}$ থেকে ক্ষুদ্রতম প্রথম স্থানাঙ্ক নেই এমন
যুগলগুলি বাদ দিই। অবশিষ্টগুলির মধ্যে ক্ষুদ্রতম
দ্বিতীয় স্থানাঙ্কবিশিষ্ট যুগলটি $X$-এর
$\canonord$-ক্ষুদ্রতম সদস্য।
\end{proof}
\noindent
এবার একটি খুব ছোট, সহজ পর্যবেক্ষণ:

\begin{prop}\ollabel{simplecardproduct}
$\cardeq{\alpha}{\beta}$ হলে
$\cardeq{\alpha\times\alpha}{\beta\times\beta}$।
\end{prop}

\begin{proof}
$f\colon\alpha\to\beta$ থেকে
$\tuple{\gamma_1,\gamma_2}\mapsto
\tuple{f(\gamma_1),f(\gamma_2)}$ চিত্রণটি পাওয়া যায়।
\end{proof}
\noindent
এগুলি কাজে লাগিয়ে একটি গুরুত্বপূর্ণ সহায়ক উপপাদ্য প্রমাণ করব:
\begin{lem}\ollabel{alphatimesalpha}
যেকোনো অসীম ক্রমসংখ্যা $\alpha$-র জন্য
$\cardeq{\alpha}{\alpha\times\alpha}$।
\end{lem}

\begin{proof}
বিরোধের জন্য ধরি, ফলটি সত্য নয় এমন ক্ষুদ্রতম
অসীম ক্রমসংখ্যা $\alpha$ আছে।
\olref[sfr][siz][zigzag]{natsquaredenumerable}
অনুযায়ী $\cardeq{\omega}{\omega\times\omega}$,
তাই $\omega\in\alpha$। তদুপরি $\alpha$ একটি
অঙ্কবাচক সংখ্যা। তা না হলে $\card{\alpha}\in\alpha$;
তাই ন্যূনতমত্বের অনুমানে
$\cardeq{\card{\alpha}}{\card{\alpha}\times\card{\alpha}}$।
আবার সংজ্ঞা অনুযায়ী $\cardeq{\card{\alpha}}{\alpha}$;
ফলে \olref{simplecardproduct} থেকে
$\cardeq{\alpha}{\alpha\times\alpha}$,
যা বিরোধ।

এবার প্রত্যেক
$\tuple{\gamma_1,\gamma_2}\in\alpha\times\alpha$-র
জন্য নিচের প্রারম্ভিক খণ্ডটি ভাবুন:
\begin{align*}
	\text{Seg}(\gamma_1, \gamma_2) &= \Setabs{\tuple{\delta_1, \delta_2} \in \alpha \times \alpha}{\tuple{\delta_1, \delta_2} \canonord \tuple{\gamma_1, \gamma_2}}
\end{align*}
$\gamma=\max(\gamma_1,\gamma_2)$ নিলে
$\tuple{\gamma_1,\gamma_2}\canonord
\tuple{\gamma+1,\gamma+1}$। তাই $\gamma$
অসীম হলে:
\begin{align*}
	\text{Seg}(\gamma_1, \gamma_2) &
	\precsim ((\gamma \ordplus 1)\ordtimes (\gamma \ordplus 1))\\
	&\approx (\gamma \ordtimes \gamma)
	\text{, অনুযায়ী \olref[ord-arithmetic][using-addition]{ordinfinitycharacter} এবং
	\olref{simplecardproduct}}\\
	&\approx \gamma \text{, আবেশ-অনুমানে}\\
	& \prec \alpha\text{, কারণ $\alpha$ অঙ্কবাচক সংখ্যা}
\end{align*}
% BN-SRC-455: সসীম সর্বোচ্চ স্থানাঙ্কের বাদ পড়া দফা আলাদা।
সর্বোচ্চ স্থানাঙ্কটি সসীম হলে খণ্ডটিও সসীম,
তাই সেক্ষেত্রেও মূল অসীম ক্রমসংখ্যার চেয়ে ছোট। অতএব
$\ordtype{\alpha\times\alpha,\canonord}\leq\alpha$,
ফলে $\cardle{\alpha\times\alpha}{\alpha}$। অবশ্যই
$\cardle{\alpha}{\alpha\times\alpha}$;
শ্র্যোডার--বার্নস্টাইন থেকে ফলটি মেলে।
% BN-SRC-809 TARGET-CORRECTION-BEGIN
\emph{উৎস-সংশোধন-নোট।} সব যথার্থ খণ্ডের সদস্যসংখ্যা
ছোট থেকে ক্রমপ্রকারের সীমায় আরেকটি ধাপ চাই। যদি পুরো
ক্রমপ্রকার $\rho>\alpha$ হত, তবে তার ঠিক $\alpha$
স্থানটির পূর্বখণ্ডের ক্রমপ্রকার $\alpha$, তাই সদস্যসংখ্যাও
$\alpha$। কিন্তু সেই যথার্থ খণ্ডের সদস্যসংখ্যা কম বলে
ইতিমধ্যে প্রমাণিত; বিরোধ। তাই $\rho\leq\alpha$।
এখানে $\alpha$ প্রাথমিক ক্রমসংখ্যা হওয়া অপরিহার্য;
নির্বিশেষ ordinal-এ ছোট সদস্যসংখ্যা মানেই ছোট ক্রমপ্রকার নয়।
% BN-SRC-809 TARGET-CORRECTION-END
\end{proof}

শেষ পর্যন্ত সরলীকরণের ফলটি পাই:

\begin{thm}\ollabel{cardplustimesmax}
$\cardfont{a},\cardfont{b}$ অসীম অঙ্কবাচক সংখ্যা হলে:
\[
	\cardfont{a}
\cardtimes \cardfont{b} = \cardfont{a} \cardplus \cardfont{b} =
\text{max}(\cardfont{a}, \cardfont{b}).
\]
\end{thm}

\begin{proof}
সাধারণত্ব না হারিয়ে
$\cardfont{a}=\max(\cardfont{a},\cardfont{b})$
ধরি। \olref{alphatimesalpha} প্রয়োগ করলে
$\cardfont{a}\cardtimes\cardfont{a}=\cardfont{a}\leq
\cardfont{a}\cardplus\cardfont{b}\leq
\cardfont{a}\cardplus\cardfont{a}\leq
\cardfont{a}\cardtimes\cardfont{a}$।
% BN-SRC-811 TARGET-CORRECTION-BEGIN
\emph{উৎস-সংশোধন-নোট।} প্রদর্শিত চেইনে যোগের
সমতাই সরাসরি প্রতিষ্ঠিত। গুণের অংশের জন্য
$1\leq\cardfont{b}\leq\cardfont{a}$ ব্যবহার করি:
দ্বিতীয় factor-এর একটি নির্দিষ্ট সদস্যে প্রথম factor-টি
বসিয়ে এবং বড় বর্গে অন্তর্ভুক্ত করে
$\cardfont{a}\leq\cardfont{a}\cardtimes\cardfont{b}
\leq\cardfont{a}\cardtimes\cardfont{a}=\cardfont{a}$।
অতএব গুণফলও একই বৃহত্তম cardinal।
% BN-SRC-811 TARGET-CORRECTION-END
\end{proof}\noindent একইভাবে $\cardfont{a}$ অসীম হলে,
প্রতিটি সেটের আকার $\leq\cardfont{a}$ এমন
$\cardfont{a}$-সংখ্যক সেটের সংযোগের আকারও
$\leq\cardfont{a}$:

\begin{prop}\ollabel{kappaunionkappasize}
$\cardfont{a}$ একটি অসীম অঙ্কবাচক সংখ্যা হোক।
প্রত্যেক ক্রমসংখ্যা $\beta\in\cardfont{a}$-র জন্য
$X_\beta$ এমন সেট হোক, যার
$\card{X_\beta}\leq\cardfont{a}$। তাহলে
$\card{\bigcup_{\beta\in\cardfont{a}}X_\beta}
\leq\cardfont{a}$।
\end{prop}

\begin{proof}
প্রত্যেক $\beta\in\cardfont{a}$-র জন্য একটি
!!a{injection} $f_\beta\colon X_\beta\to\cardfont{a}$
স্থির করি।\footnote{এগুলি কীভাবে ``স্থির'' করা যায়?
দেখুন \olref[sth][choice][countablechoice]{sec}।}
$g\colon\bigcup_{\beta\in\cardfont{a}}X_\beta\to
\cardfont{a}\times\cardfont{a}$ এভাবে সংজ্ঞায়িত
করি: $v\in X_\beta$ এবং সব
$\gamma\in\beta$-র জন্য $v\notin X_\gamma$
হলে $g(v)=\tuple{\beta,f_\beta(v)}$। এই চিত্রণটি
একটি !!a{injection}। এখন
$\bigcup_{\beta\in\cardfont{a}}X_\beta
\preceq\cardfont{a}\times\cardfont{a}
\approx\cardfont{a}$,
\olref{cardplustimesmax} অনুযায়ী।
\end{proof}

\end{document}
